\documentclass[../DoAn.tex]{subfiles}
\begin{document}

\section{Đặt vấn đề}
\label{section:1.1}
\begin{itemize}[nosep, leftmargin=*]
    \item Bối cảnh chuyển đổi số và nhu cầu tra cứu pháp luật lao động của NLĐ, NSDLĐ, luật sư nội bộ
    \item Đặc thù văn bản pháp luật VN: phân cấp Phần $\rightarrow$ Chương $\rightarrow$ Mục $\rightarrow$ Điều $\rightarrow$ Khoản $\rightarrow$ Điểm; dẫn chiếu chéo BLLĐ $\leftrightarrow$ Nghị định $\leftrightarrow$ Thông tư
    \item Quy mô corpus: 581 Điều, 2069 Khoản trên 7 văn bản pháp luật lao động
    \item Thách thức câu hỏi single-hop, multi-hop và kiểm tra điều khoản HĐLĐ
    \item Hạn chế chatbot LLM thuần: hallucination, trích sai Điều, không kiểm soát hiệu lực văn bản
\end{itemize}

\section{Các giải pháp hiện tại và hạn chế}
\label{section:1.2}

\subsection{Phương pháp tra cứu pháp luật truyền thống}
\label{subsection:1.2.1}
\begin{itemize}[nosep, leftmargin=*]
    \item Tra cứu từ khóa trên cổng thông tin điện tử, thư viện pháp luật
    \item BM25/TF-IDF, Elasticsearch: mạnh keyword, yếu ngữ nghĩa và quan hệ chéo giữa văn bản
    \item Không tổng hợp, không suy luận chuỗi quy phạm $\rightarrow$ chế tài
\end{itemize}

\subsection{Phương pháp dựa trên LLM và RAG truyền thống}
\label{subsection:1.2.2}
\begin{itemize}[nosep, leftmargin=*]
    \item ChatGPT/Gemini tra cứu chung: thiếu corpus đã index, rủi ro điều luật lỗi thời
    \item Naive RAG (chunk + embedding): cắt giữa Khoản, mất ngữ cảnh pháp lý
    \item Vector search không nắm quan hệ \texttt{obligates}, \texttt{penalizes}, \texttt{cites}
    \item Khó trả lời câu hỏi tổng hợp (ví dụ: toàn bộ quyền NLĐ khi nghỉ việc)
\end{itemize}

\subsection{Phương pháp dựa trên đồ thị tri thức và GraphRAG}
\label{subsection:1.2.3}
\begin{itemize}[nosep, leftmargin=*]
    \item Knowledge Graph thủ công: chính xác nhưng tốn công, khó mở rộng khi văn bản sửa đổi
    \item Microsoft GraphRAG: tự động trích entity/relationship, community detection, Local + Global Search
    \item Hạn chế khi áp dụng trực tiếp: schema mặc định không phù hợp domain luật lao động; thiếu lớp cấu trúc văn bản và kiểm chứng quy tắc HĐLĐ
\end{itemize}

\section{Mục tiêu và định hướng giải pháp}
\label{section:1.3}

\subsection{Mục tiêu của đồ án}
\label{subsection:1.3.1}
\begin{itemize}[nosep, leftmargin=*]
    \item \textbf{Tổng quát:} Xây dựng Trợ lý Ảo Luật Lao Động VN dựa trên GraphRAG, hỗ trợ hỏi đáp có trích dẫn và phân tích HĐLĐ
    \item Thiết kế ontology hai lớp cho corpus luật lao động
    \item Xây dựng pipeline chuẩn bị dữ liệu deterministic (Lớp 1) + indexing GraphRAG (Lớp 2)
    \item Triển khai Local Search, Global Search, multi-hop, temporal filter, rule validator
    \item Phát triển API FastAPI + frontend Next.js; module phân tích HĐLĐ (VR001--VR016)
    \item Đánh giá trên bộ benchmark \texttt{evaluation\_suite.py}
\end{itemize}

\subsection{Định hướng giải pháp}
\label{subsection:1.3.2}
\begin{itemize}[nosep, leftmargin=*]
    \item Kết hợp truy xuất có cấu trúc (đồ thị pháp luật) và truy xuất ngữ nghĩa (embedding)
    \item Lớp 1 build bằng rule (không LLM) $\rightarrow$ cấu trúc Điều/Khoản chính xác
    \item Lớp 2 extract bằng LLM với prompt domain tiếng Việt
    \item Merge graph $\rightarrow$ multi-hop trên quan hệ \texttt{entitles}, \texttt{penalizes}, \texttt{cites}, \ldots
    \item Kiểm soát đầu ra: trích dẫn Điều, lọc văn bản hết hiệu lực, rules cứng cho HĐLĐ
\end{itemize}

\section{Đóng góp chính}
\label{section:1.4}
\begin{itemize}[nosep, leftmargin=*]
    \item Ontology luật lao động hai lớp (\texttt{settings.yaml}, \texttt{02\_ontology\_design.md})
    \item Chunking theo Điều, metadata \texttt{norm\_type} (\texttt{01\_prepare\_data.py})
    \item Prompt extraction + community report tiếng Việt
    \item Lớp query nâng cao: multi-hop, temporal filter, rule validator
    \item Ứng dụng thực tế: API, frontend, contract analysis
    \item Bộ đánh giá \texttt{evaluation\_suite.py}
\end{itemize}

\end{document}
